% ID-5-02 (TO) · Mathematical Foundations of Diffusion Models · IIT Jammu
% HW 3 — From a Known Score to Langevin Sampling
% HW3 release handout. Build: xelatex hw3.tex
\documentclass[11pt,a4paper]{article}

\usepackage[margin=2.4cm]{geometry}
\usepackage{amsmath,amssymb}
\usepackage{booktabs}
\usepackage[colorlinks=true,linkcolor=black,urlcolor=blue]{hyperref}

\newcommand{\pts}[1]{\hfill\textbf{(#1pts)}}
\newcommand{\pt}[1]{\hfill\textbf{(#1pt)}}
\newcommand{\total}[1]{\par\medskip\noindent\textbf{Total: #1 points}}
\newcommand{\PP}{\mathbb{P}}
\newcommand{\EE}{\mathbb{E}}
\newcommand{\Var}{\operatorname{Var}}
\newcommand{\sig}{\operatorname{\sigma}}

% (a) (b) (c) labelling, without enumitem


\newenvironment{parts}
  {\begin{enumerate}\setlength{\itemsep}{0.45em}\setlength{\parskip}{0pt}}
  {\end{enumerate}}


\usepackage{iftex}
\ifPDFTeX\usepackage[T1]{fontenc}\usepackage[utf8]{inputenc}\fi
\providecommand{\tightlist}{\setlength{\itemsep}{0pt}\setlength{\parskip}{0pt}}
\setcounter{secnumdepth}{-1}
\setlength{\emergencystretch}{2em}
\begin{document}
\begin{center}
  {\Large\bfseries HW 3 — From a Known Score to Langevin Sampling}\\[0.4em]
  {\small Mathematical Foundations of Diffusion Models · Monsoon 2026 · IIT Jammu}
\end{center}
\medskip
\begin{center}
\begin{tabular}{@{}ll@{}}
\toprule
\textbf{Released} & Friday 23 October 2026 \\
\textbf{Due} & Monday 26 October 2026 \\
\textbf{Marks} & 40 points --- graded on correctness \\
\bottomrule
\end{tabular}
\end{center}
\medskip
\noindent\textbf{Individual homework: derivations and programming.}
Show your working and explain your numerical results. Submit an executed notebook
with the mathematical steps, code, labeled figures, and comparison table.
\bigskip\hrule\bigskip
\subsection{Background: the sampling
problem}\label{background-the-sampling-problem}

Suppose we know a probability density but want to generate observations
from it using its score. In this homework, the target is the
\textbf{standard logistic distribution} on the real line:

\[
p(x)=\frac{e^{-x}}{(1+e^{-x})^2},\qquad
F(x)=\frac{1}{1+e^{-x}},\qquad x\in\mathbb R.
\]

This is a smooth, non-Gaussian density with mean zero and variance
\(\pi^2/3\). We will derive its score, use it to construct a Langevin
sampler, and compare the generated samples with the known target.

Work entirely in \textbf{one dimension}. Write

\[s(x)=\frac{d}{dx}\log p(x).\]

The time variable \(\tau\) denotes elapsed sampling time. The target
\(p\) stays fixed as \(\tau\) increases. This differs from
diffusion-model time, which indexes a changing family of noisy data
distributions.

The distinction between a deterministic velocity and a stochastic drift
matters here. We will first examine the ODE \(\dot x=s(x)\), then use
the same score as the drift of an SDE with an additional noise term. You
will show why the latter can preserve the target distribution while the
former cannot.

\textbf{How to read this assignment.} ``Background'' and ``Given''
blocks supply explanations, formulas, or assumptions; they are not
additional questions. The lettered \textbf{Tasks}, the \textbf{Required
implementation and experiment settings}, and the \textbf{Submission}
section specify required work. The marks total 40.

\subsection{Problem 1 --- Derive and interpret the
score}\label{problem-1-derive-and-interpret-the-score}

\subsubsection{(a) Differentiate the log density (4
points)}\label{a-differentiate-the-log-density-4-points}

Starting from \(p(x)\), derive \(s(x)\) with all intermediate steps.
Express your answer both in terms of \(F(x)\) and in terms of the
hyperbolic tangent.

\par\noindent\textbf{Given: hyperbolic tangent identity}\par\nopagebreak[4]
\[\tanh(u)=\frac{e^u-e^{-u}}{e^u+e^{-u}}.\]

\subsubsection{(b) Interpret the vector field (4
points)}\label{b-interpret-the-vector-field-4-points}

Plot \(p(x)\) and \(s(x)\) on \([-8,8]\) in separate panels. Determine
the sign of \(s(x)\) for positive and negative \(x\), locate its zero,
and calculate its limits as \(x\to\pm\infty\). Explain how the score
directs the motion of a particle at each position.

\noindent\textbf{Total: 8 points}

\bigskip\hrule\bigskip

\subsection{Problem 2 --- From deterministic motion to a stochastic
sampler}\label{problem-2-from-deterministic-motion-to-a-stochastic-sampler}

\subsubsection{Given: dynamics and numerical
rules}\label{given-dynamics-and-numerical-rules}

The deterministic score ODE is

\[\frac{dx}{d\tau}=s(x),\qquad x(0)=x_{\mathrm{init}}.\]

Its explicit Euler update is

\[x_{k+1}=x_k+\eta s(x_k).\]

The Langevin SDE adds Brownian noise:

\[dX_\tau=s(X_\tau)\,d\tau+\sqrt{2}\,dW_\tau.\]

Its \textbf{Euler--Maruyama update is provided}:

\[x_{k+1}=x_k+\eta s(x_k)+\sqrt{2\eta}\,z_k,
\qquad z_k\overset{\mathrm{iid}}{\sim}\mathcal N(0,1).\]

A Brownian increment over a step of length \(\eta\) has variance
\(\eta\). Multiplication by the SDE coefficient \(\sqrt{2}\) therefore
gives the numerical noise amplitude \(\sqrt{2\eta}\). \textbf{Neither
numerical rule needs to be derived. No stochastic-calculus derivation is
required.}

\par\noindent\textbf{Background: fixed score, initialization, and noise (no proof
required)}\par\nopagebreak[4]
The score \(s(x)=\partial_x\log p(x)\) has \textbf{no explicit
dependence on sampling time} \(\tau\) or step index \(k\). Every
Euler--Maruyama step evaluates the same function at the current position
\(x_k\). Its value can change because the particle moves; the target
density and its score field remain fixed. Changing the numerical step
size does not change the target score.

For the logistic target in this homework, \(p(x)>0\) for every
\(x\in\mathbb R\), and its score is smooth, bounded, and globally
Lipschitz. Thus the Langevin SDE is well-defined from every finite
starting position. We may also draw \(X_0\sim P_{\mathrm{init}}\) from
an initial probability distribution on \(\mathbb R\), independently of
the driving Brownian motion. Examples include a Gaussian, a uniform
distribution on a finite interval, or a point mass at zero or six.

\textbf{The initial law need not have the same support as the target.}
Starting all particles at zero is valid even though the logistic target
has support throughout \(\mathbb R\). Brownian noise allows the
particles to spread into regions absent from the initial distribution;
the score drift alone would keep particles initialized at zero there.
Having full support, however, does not mean having the correct
probabilities: reaching the target distribution still requires
convergence.

This freedom of initialization is not a claim about every target
density. If \(p\) vanishes on a region, the expression
\(\partial_x\log p\) is generally undefined there, so we cannot simply
initialize there and expect the drift to repair the problem. Targets
confined to a domain may need boundary conditions, such as reflection.
Disconnected supports may prevent movement between components. For
general targets, the SDE must be well-defined at the starting points,
and convergence requires appropriate accessibility and ergodicity
conditions. Equality of initial and target supports is neither required
in general nor sufficient to establish convergence.

Initialization is separate from the fresh \(z_k\sim\mathcal N(0,1)\)
used at every numerical step. Each \(z_k\) is independent of \(X_0\) and
all earlier increments, and therefore of the current state \(x_k\). The
initial law may itself be Gaussian, but its draw and the subsequent
Gaussian increments have different roles. The score term is the
\textbf{drift}; the Brownian term is the \textbf{diffusion or noise
term}, not a Brownian drift.

For this logistic target, the exact continuous-time Langevin process
converges to the target from the stated initializations; you may use
this fact without proving it. Initialization affects how long
convergence takes. This does not imply exact target samples at a finite
running time, and the Euler--Maruyama implementation introduces an
additional discretization error. Problem 4 distinguishes these issues
from stationarity and finite-sample fluctuations.

\subsubsection{(a) Examine deterministic score motion (4
points)}\label{a-examine-deterministic-score-motion-4-points}

\begin{enumerate}
\def\labelenumi{\arabic{enumi}.}
\tightlist
\item
  Show that every trajectory of the score ODE approaches zero as
  \(\tau\to\infty\). A sign and monotonicity argument is sufficient; a
  closed-form solution is not required.
\item
  Explain why this limiting behavior does not reproduce the logistic
  target density.
\item
  Plot explicit Euler trajectories from
  \(x_{\mathrm{init}}\in\{-6,-3,-1,1,3,6\}\) with step size \(0.02\)
  until \(\tau=20\).
\end{enumerate}

\subsubsection{(b) Substitute the score into the Langevin sampler (6
points)}\label{b-substitute-the-score-into-the-langevin-sampler-6-points}

\begin{enumerate}
\def\labelenumi{\arabic{enumi}.}
\tightlist
\item
  Substitute your Problem 1 score into the given Langevin SDE and
  Euler--Maruyama update. Write the update explicitly in terms of
  \(x_k\), \(\eta\), and \(z_k\), and identify its drift and noise
  terms.
\item
  For \(\eta=0.02\), explain why the update without noise moves nonzero
  particles toward zero. State what happens if all particles start
  exactly at zero.
\item
  Explain why restoring independent Gaussian increments allows particles
  to spread from this common initial position, whereas the deterministic
  score drift alone cannot maintain the logistic target's nonzero
  variance.
\end{enumerate}

\noindent\textbf{Total: 10 points}

\bigskip\hrule\bigskip

\subsection{Problem 3 --- Implement the
sampler}\label{problem-3-implement-the-sampler}

\subsubsection{Required implementation and experiment
settings}\label{required-implementation-and-experiment-settings}

Implement the given Euler--Maruyama update, with your explicit logistic
score substituted in Problem 2(b), using NumPy or an equivalent
numerical library. At each step, evaluate the score at the current
particle positions, draw fresh standard Gaussian values, and apply the
update. Write this loop yourself; you do not need to derive a numerical
solver. No neural network is needed because the score is available
analytically.

Use the following experiment:

\begin{itemize}
\tightlist
\item
  Simulate \(M=5{,}000\) independent particles, initially all at zero.
\item
  Compare step sizes \(\eta\in\{0.2,0.05,0.01\}\).
\item
  Run each experiment to the same elapsed time \(\tau_{\max}=40\), using
  \(K=\tau_{\max}/\eta\) steps.
\item
  Draw fresh, independent Gaussian noise for every particle at every
  step.
\item
  Retain snapshots at \(\tau=0,2,10,40\) and record your random seed.
\end{itemize}

Use only the \textbf{final position of each independent particle} for
the main distribution comparisons.

\par\noindent\textbf{Background: independent observations}\par\nopagebreak[4]
Successive positions along a single trajectory are correlated. Using one
final position from each independently simulated particle avoids
treating those successive positions as independent observations.

\subsubsection{(a) Implement and show the evolving distribution (5
points)}\label{a-implement-and-show-the-evolving-distribution-5-points}

For \(\eta=0.05\), plot the particle histograms at the positive snapshot
times and overlay the target density. Describe how the initially
concentrated population spreads and approaches the target.

\subsubsection{(b) Compare numerical resolutions (5
points)}\label{b-compare-numerical-resolutions-5-points}

For each step size, report the final sample mean, variance, and elapsed
computational time. Repeat the \(\eta=0.05\) experiment with all
particles initially at \(x=6\) and compare the final empirical CDFs from
the two initializations. Discuss whether these runs provide evidence
that initialization effects have become small.

\noindent\textbf{Total: 10 points}

\bigskip\hrule\bigskip

\subsection{Problem 4 --- Explain and check the target
distribution}\label{problem-4-explain-and-check-the-target-distribution}

\subsubsection{Given: density evolution and comparison
formulas}\label{given-density-evolution-and-comparison-formulas}

Let \(q_\tau(x)\) denote the evolving density of the Langevin particles
and \(p(x)\) the fixed target. The Fokker--Planck equation is provided:

\[\frac{\partial q_\tau}{\partial\tau}
=-\frac{\partial}{\partial x}\bigl[s(x)q_\tau(x)\bigr]
+\frac{\partial^2q_\tau(x)}{\partial x^2}.\]

The empirical CDF error is

\[D_M=\sup_x\left|\widehat F_M(x)-F(x)\right|.\]

For sorted samples \(x_{(1)}\le\cdots\le x_{(M)}\), use the supplied
formula

\[D_M=\max_{1\le i\le M}\left\{
\frac{i}{M}-F(x_{(i)}),\quad
F(x_{(i)})-\frac{i-1}{M}\right\}.\]

An independent exact logistic observation can be generated by

\[U\sim\operatorname{Uniform}(0,1),\qquad X=\log\frac{U}{1-U}.\]

No derivation of these three formulas is required. The inverse-CDF
sampler is supplied only to construct a reference sample, not to replace
the Langevin implementation.

\par\noindent\textbf{Background: what the diagnostics tell us}\par\nopagebreak[4]
\textbf{Density plots are qualitative diagnostics; empirical CDF errors
quantify discrepancies.} A histogram estimates the generated density,
and its appearance depends on sample size and bin choices. \(D_M\)
compares the empirical CDF with the known target without choosing bins.
Here the target plays the role of the data distribution, but there is no
training stage because its score is known analytically.

Even exact independent target samples have nonzero \(D_M\) due to
finite-sample fluctuations. Repeated exact samples give a reference for
interpreting whether the Langevin discrepancy is small or large at the
same sample size. The range from 20 reference runs is descriptive, not a
formal acceptance threshold.

Finite running time, numerical discretization, and finite-sample
fluctuations are distinct sources of discrepancy. Smaller step sizes
reduce discretization error but need not produce a smaller observed
\(D_M\) in every experiment. Plots and sample comparisons provide
evidence of approximation, not a proof of exact equality for a
finite-step sampler.

\par\noindent\textbf{Stationarity and convergence are different statements --- not an
assessed proof}\par\nopagebreak[4]
For this logistic Langevin SDE, the target density is stationary: if the
particles already follow the target distribution, their distribution
remains the target as time passes. The score drift concentrates
probability toward the center, while Brownian noise spreads it; at the
target density, these effects balance.

Stationarity alone does \textbf{not} show that particles initialized
from a different distribution converge to the target. Convergence
requires a separate argument. For this homework, convergence of the
exact logistic Langevin process from the stated initializations is
supplied as a fact. \textbf{You are not asked to prove convergence,
establish ergodicity, or derive conditions guaranteeing either.} The
simulations provide a numerical comparison at finite running times, not
a convergence proof.

\subsubsection{(a) Verify stationarity (6
points)}\label{a-verify-stationarity-6-points}

\begin{enumerate}
\def\labelenumi{\arabic{enumi}.}
\tightlist
\item
  Rewrite the supplied Fokker--Planck equation as
  \(\partial_\tau q_\tau=-\partial_x J_\tau\) and identify the
  probability flux \(J_\tau\).
\item
  Show that \(s(x)p(x)=p'(x)\) and hence that the flux is zero when
  \(q_\tau=p\). Conclude that the target is stationary for the
  continuous-time Langevin SDE.
\item
  State the conclusion precisely: if \(q_0=p\), then \(q_\tau=p\) for
  all sampling times \(\tau\ge0\). No convergence argument is required.
\end{enumerate}

\subsubsection{(b) Compare generated samples with the target (6
points)}\label{b-compare-generated-samples-with-the-target-6-points}

\begin{enumerate}
\def\labelenumi{\arabic{enumi}.}
\tightlist
\item
  \textbf{Plot densities.} For each step size, overlay the analytic
  target density \(p(x)\) and a density-normalized histogram of the
  final generated samples. Label them ``Target density'' and ``Generated
  samples.'' Use the same bin edges and axis limits across panels.
  Describe differences in the center, spread, and tails.
\item
  \textbf{Quantify discrepancies.} Plot each final empirical CDF against
  \(F\), compute \(D_M\) using the supplied formula, and add it to the
  comparison table from Problem 3.
\item
  \textbf{Construct an exact-sampling reference.} Generate \(M=5{,}000\)
  independent logistic observations using the supplied inverse-CDF
  formula. In a reference panel, overlay their density-normalized
  histogram with \(p\), using the same bins and axes as above. Compute
  their \(D_M\) against \(F\).
\item
  \textbf{Calibrate the error size.} Repeat the exact-sampling
  experiment 20 times with independent draws. Report the median and
  range of its \(D_M\) values alongside the Langevin errors. State
  whether the Langevin discrepancies are comparable to ordinary
  finite-sample fluctuations or noticeably larger.
\item
  \textbf{Interpret the results.} Discuss finite running time,
  discretization error, and finite-sample fluctuations separately.
  Explain what your numerical results support and what they do not
  establish.
\end{enumerate}

\noindent\textbf{Total: 12 points}

\bigskip\hrule\bigskip

\subsection{Submission}\label{submission}

Submit one executed Jupyter notebook containing the derivations in
LaTeX, code, labeled figures, comparison table, and conclusions. Include
package versions and seeds. The notebook should run on a laptop CPU and
explain each experiment before presenting its code and output.

\par\noindent\textbf{Background: scope of the conclusion}\par\nopagebreak[4]
Keep the main conclusion precise: \textbf{the logistic density is
invariant for the continuous-time Langevin SDE; a sufficiently long,
sufficiently well-resolved Euler--Maruyama simulation approximates that
density.} The deterministic score ODE has a different limiting behavior.

\end{document}
